% !TeX root = ../深度学习笔记.tex
% 由主文档通过 \input 载入；请编译主文档。
\chapter*{学习检查与进阶阅读}
\addcontentsline{toc}{chapter}{学习检查与进阶阅读}
基础阶段应能独立写出数据到损失的计算过程，核对网络各层形状，推导小模型梯度并完成一次参数更新；架构阶段应能解释卷积、状态递推与注意力各自允许哪些信息交互；进阶阶段应能区分表示目标、概率模型、采样算法和评价指标。对于任何新方法，先确定改变的是模型、目标、数据还是求解算法，再比较其假设与代价，这比记忆名称更有迁移价值。

本书给出紧凑的本科入门主线，证明限于文中明确陈述并推导的结论。通用逼近、线性低秩逼近等引用结论和连续时间生成等进阶方向不构成完整专题论述。数学与基础网络可结合文献\cite{goodfellow2016}深入阅读，实现练习可结合文献\cite{d2l}；进一步研究应按需要补充统计学习理论、数值优化、计算机视觉、自然语言处理或随机过程知识。章末实验是可执行的学习任务，文中没有把未经运行的预期写成实验结果。

